\section*{الفصل \ref{ch:g5:irreversible-changes} --- تغيّرات لا رجعة فيها}

\begin{solution}{exo:g5:irreversible-changes:1}
انصهار الشوكولاتة: يمكن الرجوع عنه (تتصلب من جديد حين تبرد). \omterm{def:g4:what-a-fire-needs:burning}{احتراق} عود
الثقاب: لا يمكن. \omterm{def:g2:mixing-and-dissolving:dissolve}{ذوبان} الملح: يمكن الرجوع عنه (\omterm{def:g3:separating-mixtures:evaporate}{بتبخير} الماء).
تحميص الخبز: لا يمكن.
\end{solution}

\begin{solution}{exo:g5:irreversible-changes:2}
تغير تظهر فيه مواد جديدة. في المطبخ: طهو البيضة، وخَبز الكعكة (وأيضًا:
تحميص الخبز، واسمرار التفاحة المقطوعة).
\end{solution}

\begin{solution}{exo:g5:irreversible-changes:3}
\omterm{def:g3:separating-mixtures:evaporate}{بتبخير} الماء: يبقى الملح في الصحن. \omterm{def:g2:mixing-and-dissolving:dissolve}{فالذوبان} \omterm{def:g5:irreversible-changes:reversible}{تغير
عكوس}.
\end{solution}

\begin{solution}{exo:g5:irreversible-changes:4}
رائحة جديدة، وظهور كتل (جسم صلب) في الحليب.
\end{solution}

\begin{solution}{exo:g5:irreversible-changes:5}
التغيرات التي على اليسار: \omterm{def:g2:mixing-and-dissolving:dissolve}{ذوبان} السكر في الماء، وانصهار الجليد، \omterm{def:g2:mixing-and-dissolving:mix}{وخلط}
الرمل بالحصى. والسهم المزدوج يعني أن التغير يمكن أن يسير في الاتجاهين:
أي يمكن الرجوع عنه.
\end{solution}

\begin{solution}{exo:g5:irreversible-changes:6}
الزيت يُبعد \omterm{def:g4:air-a-mixture-of-gases:air}{الهواء} والماء عن حديد السلسلة: فهو يبطئ
الصدأ.
\end{solution}

\begin{solution}{exo:g5:irreversible-changes:7}
ظهور فقاعات غاز في سائل دون أي تسخين.
\end{solution}

\begin{solution}{exo:g5:irreversible-changes:8}
نافع: طهو الطعام، وخَبز الخبز، وحرق الغاز للتدفئة. ضار: الصدأ، وفساد
الطعام، وتعفن الخشب.
\end{solution}

\begin{solution}{exo:g5:irreversible-changes:9}
البرد يبطئ التغيرات الكيميائية التي تجعل الحليب يتحمض، فيحفظ مدة
أطول.
\end{solution}

\begin{solution}{exo:g5:irreversible-changes:10}
لا. فالفقاعات بخار ماء يعود ماءً حين يبرد: لا يُصنع شيء جديد، ويمكن
الرجوع عن التغير. والفقاعات وحدها لا تثبت أن مادة جديدة قد
ظهرت.
\end{solution}

\begin{solution}{exo:g5:irreversible-changes:11}
الشمعة المشتعلة تطلق حرارة وضوءًا، وتصنع مواد جديدة: بخار الماء (تتغشى
كأس باردة فوقها) وثنائي أكسيد الكربون. والشمع الذي يحترق لا يعود. أما
الشمع المنصهر فيتصلب من جديد حين يبرد.
\end{solution}
